% Tabla A1-T1. Generada automáticamente; no editar a mano.
% Requiere \usepackage{booktabs} (y graphicx si la tabla es ancha).
\begin{table}[htbp]
    \centering
    \caption{Análisis I: composición de la muestra de planetas en tránsito por régimen de radio.}
    \label{tab:A1-T1}
    \small
    \resizebox{\textwidth}{!}{%
    \begin{tabular}{llrrrrr}
    \toprule
    Régimen & Criterio & $N_{\rm pl}$ & \% & $P$ med. [d] & $M_\star$ med. [$M_\odot$] & {[Fe/H]} med. \\
    \midrule
    Terrestre & Rp $<$ 1.25 $R_\oplus$ & 514 & 11.5 & 4.46 & 0.860 & -0.010 \\
    Super-Tierra & 1.25 $\leq$ Rp $<$ 2 $R_\oplus$ & 1\,080 & 24.2 & 6.50 & 0.910 & 0.000 \\
    Tipo Neptuno & 2 $\leq$ Rp $<$ 6 $R_\oplus$ & 1\,937 & 43.4 & 14.55 & 0.939 & 0.010 \\
    Gigante gaseoso & 6 $\leq$ Rp $<$ 15 $R_\oplus$ & 720 & 16.1 & 4.62 & 1.034 & 0.120 \\
    Mayor que Júpiter & Rp $\geq$ 15 $R_\oplus$ & 210 & 4.7 & 3.28 & 1.322 & 0.064 \\
    \bottomrule
    \end{tabular}
    }
    \par\smallskip
    {\footnotesize\textit{Nota:} N = 4461 planetas alrededor de 3390 estrellas.}
\end{table}
